% =====================================================================
%  state_layer.tex  --  PHASE II, Tasks 3 and 4 (+ 4b, 4c)
%
%  Results for the state-inference stage: what separates STANDBY, how
%  the decoding rule compares against clustering baselines, why a
%  transition prior cannot remove flicker, and how a partition is
%  validated with no ground truth.
%
%  Intended placement: Section 4 (\label{sec:states}), immediately after
%  the method.  Also carries \label{sec:results} (it is the first
%  results section) and \label{sec:validation}.
%
%  Drop-in usage: % =====================================================================
%  state_layer.tex  --  PHASE II, Tasks 3 and 4 (+ 4b, 4c)
%
%  Results for the state-inference stage: what separates STANDBY, how
%  the decoding rule compares against clustering baselines, why a
%  transition prior cannot remove flicker, and how a partition is
%  validated with no ground truth.
%
%  Intended placement: Section 4 (\label{sec:states}), immediately after
%  the method.  Also carries \label{sec:results} (it is the first
%  results section) and \label{sec:validation}.
%
%  Drop-in usage: % =====================================================================
%  state_layer.tex  --  PHASE II, Tasks 3 and 4 (+ 4b, 4c)
%
%  Results for the state-inference stage: what separates STANDBY, how
%  the decoding rule compares against clustering baselines, why a
%  transition prior cannot remove flicker, and how a partition is
%  validated with no ground truth.
%
%  Intended placement: Section 4 (\label{sec:states}), immediately after
%  the method.  Also carries \label{sec:results} (it is the first
%  results section) and \label{sec:validation}.
%
%  Drop-in usage: % =====================================================================
%  state_layer.tex  --  PHASE II, Tasks 3 and 4 (+ 4b, 4c)
%
%  Results for the state-inference stage: what separates STANDBY, how
%  the decoding rule compares against clustering baselines, why a
%  transition prior cannot remove flicker, and how a partition is
%  validated with no ground truth.
%
%  Intended placement: Section 4 (\label{sec:states}), immediately after
%  the method.  Also carries \label{sec:results} (it is the first
%  results section) and \label{sec:validation}.
%
%  Drop-in usage: \input{state_layer}
%  Figures are referenced by bare name; add to the preamble
%      \graphicspath{{../outputs/phase2/figures/pelletizer-I/}}
%  Bibliography: paper/references.bib
%  Numbers below are generated by
%      python run_phase2.py
%  and stored in outputs/phase2/task3/<machine>/task3_results.json,
%  outputs/phase2/task4/<machine>/{comparison_table.csv,
%  flicker_diagnosis.json,sticky_prior_sweep.csv,min_dwell_decode.csv}.
% =====================================================================

\section{The State Layer}
\label{sec:states}
\label{sec:results}

This section reports what the state-inference stage of
Section~\ref{sec:method:labelling} produces on pelletizer~I: which
measured quantities separate the states, how the choice of clustering
algorithm compares with the choice of decoding rule, and how a partition
is validated when there is nothing to validate it against.

\subsection{What separates STANDBY}
\label{sec:states:features}

Twenty-four features are derived from the six measured channels in four
groups: raw electrical (5), derived electrical (4 --- power factor,
$Q/P$ ratio, load factor, apparent-power residual), rolling statistics
(8, over a $300$\,s window), and temporal (7). Rolling windows are
\emph{trailing}, not centred, and are computed within contiguous
segments so that none spans an acquisition gap.

\textbf{The target is STANDBY against $\{$WORKING, PEAK\_LOAD$\}$ on
energised rows only.} Separating OFF from everything else is trivial and
would dominate any importance ranking computed over all four classes;
separating an energised idle machine from a lightly loaded running one
is the problem the framework actually has to solve. A four-class
ranking was also computed and is not reported here, because it measures
the OFF boundary and invites the two to be averaged.

\begin{table}[t]
\centering
\caption{Feature-group importance for the STANDBY-versus-productive
split, by three measures with different failure modes. Each column is
normalised to its own total.}
\label{tab:feature-groups}
\small
\begin{tabular}{@{}lrrr@{}}
\toprule
Group & mutual info. & split gain & SHAP \\
\midrule
raw electrical      & $0.350$ & $\mathbf{0.560}$ & $\mathbf{0.537}$ \\
derived electrical  & $0.250$ & $0.438$ & $0.310$ \\
rolling statistics  & $\mathbf{0.388}$ & $0.001$ & $0.116$ \\
temporal            & $0.012$ & $0.002$ & $0.036$ \\
\bottomrule
\end{tabular}
\end{table}

Three readings of Table~\ref{tab:feature-groups} matter for the rest of
the paper.

\textbf{The physics argument is empirically confirmed.}
\texttt{power\_factor} and \texttt{q\_p\_ratio} rank second and third by
consensus over the three measures, behind \texttt{active\_power} and
ahead of \texttt{current}, and the four derived electrical features
carry $31$--$44\%$ of total importance. Methods~1 and~2 of the
validation battery (Section~\ref{sec:validation}) were previously
justified by induction-motor theory alone; they now have a measured
basis on the machine's own data. Section~\ref{sec:explain:state}
strengthens this: under Shapley attribution power factor overtakes
active power outright.

\textbf{Temporal features are nearly irrelevant to this split}
($1$--$4\%$). That is the desirable outcome, not a null one: a state
model leaning on \texttt{hour\_of\_day} would be reproducing the shift
schedule rather than the machine's electrical condition, and would not
transfer to a facility with a different calendar. It is also the exact
opposite of what the \emph{decision} layer's forecaster does, which
draws $82\%$ of its attribution from the calendar
(Section~\ref{sec:crossmachine:decision}) --- the two stages of the
framework depend on almost disjoint information, and only the state
stage is physics-driven.

\textbf{High mutual information with near-zero split gain is the
redundancy signature.} Rolling statistics are individually informative
(MI $0.388$, the largest in the column) and add nothing once the tree
already holds the instantaneous value (gain $0.001$). Reporting either
measure alone would misdescribe them.

\paragraph{A caveat that must not be dropped.}
The surrogate used for the gain and SHAP columns reproduces its targets
to $99.9\%$ accuracy. \emph{That number is not a validation result.} The
targets are the state model's own labels and the state model was fitted
on a subset of these same features, so the surrogate is largely
re-deriving a deterministic function of its own inputs. The rankings are
meaningful; the accuracy is evidence of nothing. Validation is
Section~\ref{sec:validation}'s job.

\subsection{The emission model barely matters; the decoding rule does}
\label{sec:results:task4}

Six state-inference configurations are compared on one representation
and one protocol: $64$ contiguous two-hour blocks spread across the
record and never crossing a gap, with alternate blocks forming a fitting
set and an evaluation set of $230\,400$ rows each. DBSCAN and spectral
clustering are transductive, so their labels are extended to the
evaluation blocks by one-nearest-neighbour, applied identically to both.
Because no ground truth exists, the comparison is scored on flicker,
median dwell, the physics battery, the balanced schedule score, and the
STANDBY hour total --- never on accuracy.

\begin{table}[t]
\centering
\caption{State-inference comparison on the evaluation blocks. Flicker is
the fraction of state runs shorter than the $10$\,s physical
plausibility floor; the binary column recomputes it on the
STANDBY-versus-rest sequence the decision layer consumes. No column is
interpretable on its own --- see the spectral row.}
\label{tab:state-detection}
\small
\setlength{\tabcolsep}{4pt}
\begin{tabular}{@{}lrrrrrr@{}}
\toprule
Method & flick. & flick. & dwell & phys. & sched. & STANDBY \\
       & \%     & bin.\% & (s)   &       & \%      & (h) \\
\midrule
$k$-means \cite{lloyd1982least}      & $78.2$ & $49.1$ & $2$ & $5/5$ & $84.4$ & $4.38$ \\
DBSCAN \cite{ester1996dbscan}        & $62.5$ & $50.3$ & $4$ & $5/5$ & $84.4$ & $4.34$ \\
Spectral \cite{ng2001spectral}       & $98.2$ & $98.8$ & $2$ & $\mathbf{4/5}$ & $\mathbf{89.0}$ & $\mathbf{15.01}$ \\
GMM, i.i.d.\ \cite{mclachlan2000finite} & $66.9$ & $53.6$ & $3$ & $5/5$ & $84.5$ & $4.17$ \\
GMM--HMM \cite{rabiner1989tutorial}  & $44.5$ & $44.7$ & $13$ & $5/5$ & $84.5$ & $4.16$ \\
\textbf{GMM--HMM + min-dwell}        & $\mathbf{0.00}$ & $\mathbf{0.00}$ & $\mathbf{92.5}$ & $5/5$ & $84.5$ & $4.20$ \\
\bottomrule
\end{tabular}
\end{table}

\begin{figure*}[t]
\centering
\includegraphics[width=\textwidth]{fig_p7_state_detection}
\caption{The same six configurations as
Table~\ref{tab:state-detection}, plotted. \textbf{(A)} Flicker against
median dwell. The three filled points share one set of GMM emissions
and differ only in how those emissions are decoded; the arrows follow
i.i.d.\ assignment $\rightarrow$ Viterbi $\rightarrow$ Viterbi under a
minimum-dwell floor. The three squares are alternative emission models
decoded i.i.d.\ Changing the emission model moves the result within one
cluster; changing the decoding rule moves it two orders of magnitude in
dwell and to zero in flicker. \textbf{(B)} Inferred STANDBY hours.
Spectral clustering assigns $3.6\times$ every other method's total while
scoring the highest factory-schedule agreement of any row in the table,
which is why the battery is argued for in preference to any single
ground-truth-free metric. Hatching marks a failed physics check in both
panels.}
\label{fig:state-detection}
\end{figure*}

\textbf{Choosing a better clustering algorithm buys almost nothing;
choosing a better sequence model buys everything} ---
Figure~\ref{fig:state-detection}(A) is the table read as a picture, and
the two families separate along different axes. $k$-means, DBSCAN and
the GMM all pass $5/5$ physics checks, agree with the factory schedule
to within $0.2$ points, and find $4.17$--$4.38$ STANDBY hours on the
evaluation blocks. Switching from i.i.d.\ assignment to Viterbi decoding
over the \emph{same} emissions cuts flicker from $66.9\%$ to $44.5\%$
and quadruples the median dwell; adding the minimum-dwell constraint
takes flicker to zero and the median dwell to $92.5$\,s. The
emission-model rows span $62$--$78\%$ of flicker; the decoding rows span
$0$--$67\%$.

The comparison is clean in the one place it needs to be: the
\textsc{gmm} and \textsc{gmm--hmm} rows differ \emph{only} in the
decoding rule, because the HMM's emissions are warm-started from the
fitted GMM and held fixed (Section~\ref{sec:method:ghmm}). Fitting HMM
emissions from scratch by Baum--Welch converges to a materially worse
optimum on this record --- its ``STANDBY'' absorbs part of the
productive range at $10.3$\,kW instead of finding the no-load band at
$3.67$\,kW, and it fails the variance-ordering check --- which would
have confounded the value of temporal decoding with EM initialisation
luck.

\textbf{Spectral clustering is the one genuine failure, and it is
instructive} (Figure~\ref{fig:state-detection}(B)). It assigns $15.01$
STANDBY hours, $3.6\times$ every other
method, with $98.8\%$ binary flicker and a no-load current ratio of
$0.001$ against the acceptance band $[0.25,0.50]$. It also scores the
\emph{highest} schedule agreement in the table ($89.0\%$): a partition
that over-assigns non-productive states automatically looks good on the
factory-closed window. That is a weakness of the schedule metric in
isolation and the reason the STANDBY hour total is printed beside it.
\textbf{No single ground-truth-free metric can be quoted alone}, which
is the argument for the battery rather than for a favourite score.

\paragraph{A protocol correction worth recording.}
DBSCAN chooses its own cluster count and returned $28$. Naming the four
lowest-power clusters with the four canonical state names put all four
``states'' inside the machine's low-power range and produced an apparent
total failure --- $0.06$ STANDBY hours, $98.9\%$ flicker --- that was
entirely an artefact of the mapping. Cluster means are now grouped into
four size-weighted power bands, so sub-clusters of one physical state
are re-merged, and DBSCAN passes $5/5$ with $4.34$ STANDBY hours. Any
comparison against a method that over-segments needs this step, and
without it this paper would have reported a baseline failure that does
not exist.

\subsection{Flicker: why a transition prior cannot remove it}
\label{sec:flicker}

The expectation going in was that Viterbi decoding under a sticky prior
would take flicker to a few per cent. It does not, and the reason is
arithmetic rather than empirical.

\textbf{The sticky prior as conventionally set is numerically inert.}
Setting the Dirichlet pseudo-count to the target dwell in samples gives
$\kappa=30$ against roughly $576\,000$ observed transitions. A flat
prior and $\kappa=30$ give flicker $64.10\%$ and $64.09\%$ and minimum
self-transition $0.9677$ and $0.9678$: the prior is present in the
model, cited in its documentation, and has no measurable effect.

\textbf{Strengthening it does not help either.} Scaling the prior to ten
times the total evidence ($\kappa=575\,920$, minimum self-transition
$0.9994$, implied mean dwell $1\,605$\,s) still leaves flicker at
$50.2\%$.

\textbf{The bound is structural.} In a Viterbi decode the cost of
changing state is $\log(p_{\text{self}}/p_{\text{switch}})$, measured
here at $9.23$ nats at its largest and bounded above however strong the
prior is made. The emission log-likelihood gap between the best and
second-best state has median $27.8$ nats and $90$th percentile
$41\,701$ nats. \textbf{At $55.7\%$ of timesteps the emission term
exceeds the largest penalty the transition matrix can apply}, so the
transition model is arithmetically incapable of changing the decision
there and the HMM degenerates toward per-sample assignment. This is a
general property of applying an HMM to low-noise, multi-channel
measurements, not a property of this record.

\textbf{What fixes it is an explicit duration constraint} --- the
discrete analogue of a hidden semi-Markov duration
model~\cite{yu2010hsmm}, imposed at decode time. The floor is set equal
to the flicker threshold itself, so it removes exactly what the metric
defines as implausible and nothing more.
Table~\ref{tab:min-dwell} sweeps the floor over an order of magnitude.

\begin{table}[t]
\centering
\caption{Minimum-dwell sweep. Flicker goes to zero at no cost to the
physics checks, the schedule agreement, or the STANDBY total. This
sweep is run on the flicker-diagnosis decode rather than on the
evaluation blocks of Table~\ref{tab:state-detection}, which is why its
absolute physics score and STANDBY total differ from that table's; the
comparison of interest is down the column, not across the two tables.}
\label{tab:min-dwell}
\small
\begin{tabular}{@{}rrrrr@{}}
\toprule
min.\ dwell & flicker \% & median dwell (s) & physics & STANDBY (h) \\
\midrule
$0$\,s  & $64.09$ & $4$   & $4/5$ & $5.34$ \\
$10$\,s & $\mathbf{0.00}$ & $153$ & $4/5$ & $5.34$ \\
$30$\,s & $0.00$ & $327.5$ & $4/5$ & $5.37$ \\
$60$\,s & $0.00$ & $436.5$ & $4/5$ & $5.40$ \\
\bottomrule
\end{tabular}
\end{table}

\begin{figure}[t]
\centering
\includegraphics[width=\columnwidth]{fig_p7_min_dwell}
\caption{The minimum-dwell sweep of Table~\ref{tab:min-dwell}.
\textbf{(A)} Flicker under both definitions collapses to $0.00\%$ at the
first non-zero floor and stays there. \textbf{(B)} The median dwell
rises by two orders of magnitude while the inferred STANDBY total is
flat. The right-hand axis is anchored at zero deliberately: on an
auto-scaled axis the $1\%$ drift across the whole sweep would render as
a pronounced trend and say the opposite of what the data says.}
\label{fig:min-dwell}
\end{figure}

The constraint is removing noise, not reshaping the partition
(Figure~\ref{fig:min-dwell}): the
STANDBY total moves by $0.00$\,h from $0$\,s to $10$\,s and by $1\%$ out
to $60$\,s, and no physics check changes verdict. Its cost is measured
directly --- $0.19\%$ of rows relabelled on this machine, $0.15$--$6.1\%$
across the plant (Section~\ref{sec:multimachine:general}) --- and it is
the component the ablation of Section~\ref{sec:ablation:state} credits
with removing flicker, in preference to the HMM.

\subsection{Validation without ground truth}
\label{sec:validation}

IMDELD supplies no per-timestamp state annotation, and obtaining one
would require instrumenting each machine's contactor. The partition
therefore cannot be scored by any supervised metric, and the alternative
adopted here is to require agreement with evidence the model does not
have access to.

Three independent bodies of evidence are used, and they fail in
different ways, which is the property that makes their agreement
informative.

\textbf{(i) Circuit physics.} The five checks C1--C5 of
Section~\ref{sec:method:accept_reject}, of which C3 (power-factor
separation) and C4 (no-load current ratio) are the two that do not
depend on power magnitude and therefore the two that can distinguish an
energised idle machine from a lightly loaded one. Pelletizer~I passes
$5/5$, with a power-factor separation of $0.72$ and a no-load current
ratio of $0.36$ --- inside the $[0.25,0.50]$ band that induction-motor
theory predicts and that no part of the fitting procedure was told
about.

\textbf{(ii) Model-level diagnostics.} A ten-test harness (T1--T10)
covering data sanity, EM convergence, model order, component separation
at $2\sigma$, assignment confidence, ghost components, physical
thresholds, variance ordering, temporal alignment with nights and
weekends, and agreement with the published plant calendar. T3 is
reported as a \emph{diagnostic} rather than a selector: as
Section~\ref{sec:multimachine:general} shows, BIC does not identify $k$
on records of this size, and the physics constraints do the selecting.

\textbf{(iii) Independently derived classifiers.} Two threshold rules
constructed from the same raw channels but with no reference to the
mixture model: a power-factor classifier and a magnetising-current-ratio
classifier, each with its threshold chosen by Otsu's
criterion~\cite{otsu1979threshold} rather than by hand, so that neither
is tuned to agree. On pelletizer~I the Otsu power-factor threshold falls
at $0.474$ and the resulting classifier agrees with the state model on
$\mathbf{99.4\%}$ of energised rows; the current-ratio classifier agrees
on $97.3\%$. Agreement across the eight machines is $85$--$99.7\%$, with
one instructive exception at $55.1\%$ on the machine whose physics
checks fail hardest (Section~\ref{sec:multimachine:states}). Verdicts
are aggregated by consensus~\cite{strehl2002cluster}.

Finally, the plant's own calendar provides an external operational
check: the facility halts every weekday between 17:00 and 22:00 local
time to avoid the peak tariff, and does not run at weekends. On the full
record $76.8\%$ of readings inside that closed window are labelled
non-productive. The converse is reported alongside it, because a
degenerate partition calling everything OFF would otherwise score
$100\%$.

\textbf{The battery is not decoration, and Phase~V measured what it is
worth.} Re-fitting the labeller at ten seeds moves the headline STANDBY
total by a coefficient of variation of $33.7\%$, and three seeds in ten
converge on a cluster that is physically a different state altogether.
The battery scores exactly those three below threshold and the other
seven at $1.00$, and conditioning on a passing score collapses the
spread to $3.0\%$ (Section~\ref{sec:significance:seeds}). That is why
the accept/reject loop is part of the method rather than an appendix to
it, and why every state-time total in this paper is quoted conditional
on it.

%  Figures are referenced by bare name; add to the preamble
%      \graphicspath{{../outputs/phase2/figures/pelletizer-I/}}
%  Bibliography: paper/references.bib
%  Numbers below are generated by
%      python run_phase2.py
%  and stored in outputs/phase2/task3/<machine>/task3_results.json,
%  outputs/phase2/task4/<machine>/{comparison_table.csv,
%  flicker_diagnosis.json,sticky_prior_sweep.csv,min_dwell_decode.csv}.
% =====================================================================

\section{The State Layer}
\label{sec:states}
\label{sec:results}

This section reports what the state-inference stage of
Section~\ref{sec:method:labelling} produces on pelletizer~I: which
measured quantities separate the states, how the choice of clustering
algorithm compares with the choice of decoding rule, and how a partition
is validated when there is nothing to validate it against.

\subsection{What separates STANDBY}
\label{sec:states:features}

Twenty-four features are derived from the six measured channels in four
groups: raw electrical (5), derived electrical (4 --- power factor,
$Q/P$ ratio, load factor, apparent-power residual), rolling statistics
(8, over a $300$\,s window), and temporal (7). Rolling windows are
\emph{trailing}, not centred, and are computed within contiguous
segments so that none spans an acquisition gap.

\textbf{The target is STANDBY against $\{$WORKING, PEAK\_LOAD$\}$ on
energised rows only.} Separating OFF from everything else is trivial and
would dominate any importance ranking computed over all four classes;
separating an energised idle machine from a lightly loaded running one
is the problem the framework actually has to solve. A four-class
ranking was also computed and is not reported here, because it measures
the OFF boundary and invites the two to be averaged.

\begin{table}[t]
\centering
\caption{Feature-group importance for the STANDBY-versus-productive
split, by three measures with different failure modes. Each column is
normalised to its own total.}
\label{tab:feature-groups}
\small
\begin{tabular}{@{}lrrr@{}}
\toprule
Group & mutual info. & split gain & SHAP \\
\midrule
raw electrical      & $0.350$ & $\mathbf{0.560}$ & $\mathbf{0.537}$ \\
derived electrical  & $0.250$ & $0.438$ & $0.310$ \\
rolling statistics  & $\mathbf{0.388}$ & $0.001$ & $0.116$ \\
temporal            & $0.012$ & $0.002$ & $0.036$ \\
\bottomrule
\end{tabular}
\end{table}

Three readings of Table~\ref{tab:feature-groups} matter for the rest of
the paper.

\textbf{The physics argument is empirically confirmed.}
\texttt{power\_factor} and \texttt{q\_p\_ratio} rank second and third by
consensus over the three measures, behind \texttt{active\_power} and
ahead of \texttt{current}, and the four derived electrical features
carry $31$--$44\%$ of total importance. Methods~1 and~2 of the
validation battery (Section~\ref{sec:validation}) were previously
justified by induction-motor theory alone; they now have a measured
basis on the machine's own data. Section~\ref{sec:explain:state}
strengthens this: under Shapley attribution power factor overtakes
active power outright.

\textbf{Temporal features are nearly irrelevant to this split}
($1$--$4\%$). That is the desirable outcome, not a null one: a state
model leaning on \texttt{hour\_of\_day} would be reproducing the shift
schedule rather than the machine's electrical condition, and would not
transfer to a facility with a different calendar. It is also the exact
opposite of what the \emph{decision} layer's forecaster does, which
draws $82\%$ of its attribution from the calendar
(Section~\ref{sec:crossmachine:decision}) --- the two stages of the
framework depend on almost disjoint information, and only the state
stage is physics-driven.

\textbf{High mutual information with near-zero split gain is the
redundancy signature.} Rolling statistics are individually informative
(MI $0.388$, the largest in the column) and add nothing once the tree
already holds the instantaneous value (gain $0.001$). Reporting either
measure alone would misdescribe them.

\paragraph{A caveat that must not be dropped.}
The surrogate used for the gain and SHAP columns reproduces its targets
to $99.9\%$ accuracy. \emph{That number is not a validation result.} The
targets are the state model's own labels and the state model was fitted
on a subset of these same features, so the surrogate is largely
re-deriving a deterministic function of its own inputs. The rankings are
meaningful; the accuracy is evidence of nothing. Validation is
Section~\ref{sec:validation}'s job.

\subsection{The emission model barely matters; the decoding rule does}
\label{sec:results:task4}

Six state-inference configurations are compared on one representation
and one protocol: $64$ contiguous two-hour blocks spread across the
record and never crossing a gap, with alternate blocks forming a fitting
set and an evaluation set of $230\,400$ rows each. DBSCAN and spectral
clustering are transductive, so their labels are extended to the
evaluation blocks by one-nearest-neighbour, applied identically to both.
Because no ground truth exists, the comparison is scored on flicker,
median dwell, the physics battery, the balanced schedule score, and the
STANDBY hour total --- never on accuracy.

\begin{table}[t]
\centering
\caption{State-inference comparison on the evaluation blocks. Flicker is
the fraction of state runs shorter than the $10$\,s physical
plausibility floor; the binary column recomputes it on the
STANDBY-versus-rest sequence the decision layer consumes. No column is
interpretable on its own --- see the spectral row.}
\label{tab:state-detection}
\small
\setlength{\tabcolsep}{4pt}
\begin{tabular}{@{}lrrrrrr@{}}
\toprule
Method & flick. & flick. & dwell & phys. & sched. & STANDBY \\
       & \%     & bin.\% & (s)   &       & \%      & (h) \\
\midrule
$k$-means \cite{lloyd1982least}      & $78.2$ & $49.1$ & $2$ & $5/5$ & $84.4$ & $4.38$ \\
DBSCAN \cite{ester1996dbscan}        & $62.5$ & $50.3$ & $4$ & $5/5$ & $84.4$ & $4.34$ \\
Spectral \cite{ng2001spectral}       & $98.2$ & $98.8$ & $2$ & $\mathbf{4/5}$ & $\mathbf{89.0}$ & $\mathbf{15.01}$ \\
GMM, i.i.d.\ \cite{mclachlan2000finite} & $66.9$ & $53.6$ & $3$ & $5/5$ & $84.5$ & $4.17$ \\
GMM--HMM \cite{rabiner1989tutorial}  & $44.5$ & $44.7$ & $13$ & $5/5$ & $84.5$ & $4.16$ \\
\textbf{GMM--HMM + min-dwell}        & $\mathbf{0.00}$ & $\mathbf{0.00}$ & $\mathbf{92.5}$ & $5/5$ & $84.5$ & $4.20$ \\
\bottomrule
\end{tabular}
\end{table}

\begin{figure*}[t]
\centering
\includegraphics[width=\textwidth]{fig_p7_state_detection}
\caption{The same six configurations as
Table~\ref{tab:state-detection}, plotted. \textbf{(A)} Flicker against
median dwell. The three filled points share one set of GMM emissions
and differ only in how those emissions are decoded; the arrows follow
i.i.d.\ assignment $\rightarrow$ Viterbi $\rightarrow$ Viterbi under a
minimum-dwell floor. The three squares are alternative emission models
decoded i.i.d.\ Changing the emission model moves the result within one
cluster; changing the decoding rule moves it two orders of magnitude in
dwell and to zero in flicker. \textbf{(B)} Inferred STANDBY hours.
Spectral clustering assigns $3.6\times$ every other method's total while
scoring the highest factory-schedule agreement of any row in the table,
which is why the battery is argued for in preference to any single
ground-truth-free metric. Hatching marks a failed physics check in both
panels.}
\label{fig:state-detection}
\end{figure*}

\textbf{Choosing a better clustering algorithm buys almost nothing;
choosing a better sequence model buys everything} ---
Figure~\ref{fig:state-detection}(A) is the table read as a picture, and
the two families separate along different axes. $k$-means, DBSCAN and
the GMM all pass $5/5$ physics checks, agree with the factory schedule
to within $0.2$ points, and find $4.17$--$4.38$ STANDBY hours on the
evaluation blocks. Switching from i.i.d.\ assignment to Viterbi decoding
over the \emph{same} emissions cuts flicker from $66.9\%$ to $44.5\%$
and quadruples the median dwell; adding the minimum-dwell constraint
takes flicker to zero and the median dwell to $92.5$\,s. The
emission-model rows span $62$--$78\%$ of flicker; the decoding rows span
$0$--$67\%$.

The comparison is clean in the one place it needs to be: the
\textsc{gmm} and \textsc{gmm--hmm} rows differ \emph{only} in the
decoding rule, because the HMM's emissions are warm-started from the
fitted GMM and held fixed (Section~\ref{sec:method:ghmm}). Fitting HMM
emissions from scratch by Baum--Welch converges to a materially worse
optimum on this record --- its ``STANDBY'' absorbs part of the
productive range at $10.3$\,kW instead of finding the no-load band at
$3.67$\,kW, and it fails the variance-ordering check --- which would
have confounded the value of temporal decoding with EM initialisation
luck.

\textbf{Spectral clustering is the one genuine failure, and it is
instructive} (Figure~\ref{fig:state-detection}(B)). It assigns $15.01$
STANDBY hours, $3.6\times$ every other
method, with $98.8\%$ binary flicker and a no-load current ratio of
$0.001$ against the acceptance band $[0.25,0.50]$. It also scores the
\emph{highest} schedule agreement in the table ($89.0\%$): a partition
that over-assigns non-productive states automatically looks good on the
factory-closed window. That is a weakness of the schedule metric in
isolation and the reason the STANDBY hour total is printed beside it.
\textbf{No single ground-truth-free metric can be quoted alone}, which
is the argument for the battery rather than for a favourite score.

\paragraph{A protocol correction worth recording.}
DBSCAN chooses its own cluster count and returned $28$. Naming the four
lowest-power clusters with the four canonical state names put all four
``states'' inside the machine's low-power range and produced an apparent
total failure --- $0.06$ STANDBY hours, $98.9\%$ flicker --- that was
entirely an artefact of the mapping. Cluster means are now grouped into
four size-weighted power bands, so sub-clusters of one physical state
are re-merged, and DBSCAN passes $5/5$ with $4.34$ STANDBY hours. Any
comparison against a method that over-segments needs this step, and
without it this paper would have reported a baseline failure that does
not exist.

\subsection{Flicker: why a transition prior cannot remove it}
\label{sec:flicker}

The expectation going in was that Viterbi decoding under a sticky prior
would take flicker to a few per cent. It does not, and the reason is
arithmetic rather than empirical.

\textbf{The sticky prior as conventionally set is numerically inert.}
Setting the Dirichlet pseudo-count to the target dwell in samples gives
$\kappa=30$ against roughly $576\,000$ observed transitions. A flat
prior and $\kappa=30$ give flicker $64.10\%$ and $64.09\%$ and minimum
self-transition $0.9677$ and $0.9678$: the prior is present in the
model, cited in its documentation, and has no measurable effect.

\textbf{Strengthening it does not help either.} Scaling the prior to ten
times the total evidence ($\kappa=575\,920$, minimum self-transition
$0.9994$, implied mean dwell $1\,605$\,s) still leaves flicker at
$50.2\%$.

\textbf{The bound is structural.} In a Viterbi decode the cost of
changing state is $\log(p_{\text{self}}/p_{\text{switch}})$, measured
here at $9.23$ nats at its largest and bounded above however strong the
prior is made. The emission log-likelihood gap between the best and
second-best state has median $27.8$ nats and $90$th percentile
$41\,701$ nats. \textbf{At $55.7\%$ of timesteps the emission term
exceeds the largest penalty the transition matrix can apply}, so the
transition model is arithmetically incapable of changing the decision
there and the HMM degenerates toward per-sample assignment. This is a
general property of applying an HMM to low-noise, multi-channel
measurements, not a property of this record.

\textbf{What fixes it is an explicit duration constraint} --- the
discrete analogue of a hidden semi-Markov duration
model~\cite{yu2010hsmm}, imposed at decode time. The floor is set equal
to the flicker threshold itself, so it removes exactly what the metric
defines as implausible and nothing more.
Table~\ref{tab:min-dwell} sweeps the floor over an order of magnitude.

\begin{table}[t]
\centering
\caption{Minimum-dwell sweep. Flicker goes to zero at no cost to the
physics checks, the schedule agreement, or the STANDBY total. This
sweep is run on the flicker-diagnosis decode rather than on the
evaluation blocks of Table~\ref{tab:state-detection}, which is why its
absolute physics score and STANDBY total differ from that table's; the
comparison of interest is down the column, not across the two tables.}
\label{tab:min-dwell}
\small
\begin{tabular}{@{}rrrrr@{}}
\toprule
min.\ dwell & flicker \% & median dwell (s) & physics & STANDBY (h) \\
\midrule
$0$\,s  & $64.09$ & $4$   & $4/5$ & $5.34$ \\
$10$\,s & $\mathbf{0.00}$ & $153$ & $4/5$ & $5.34$ \\
$30$\,s & $0.00$ & $327.5$ & $4/5$ & $5.37$ \\
$60$\,s & $0.00$ & $436.5$ & $4/5$ & $5.40$ \\
\bottomrule
\end{tabular}
\end{table}

\begin{figure}[t]
\centering
\includegraphics[width=\columnwidth]{fig_p7_min_dwell}
\caption{The minimum-dwell sweep of Table~\ref{tab:min-dwell}.
\textbf{(A)} Flicker under both definitions collapses to $0.00\%$ at the
first non-zero floor and stays there. \textbf{(B)} The median dwell
rises by two orders of magnitude while the inferred STANDBY total is
flat. The right-hand axis is anchored at zero deliberately: on an
auto-scaled axis the $1\%$ drift across the whole sweep would render as
a pronounced trend and say the opposite of what the data says.}
\label{fig:min-dwell}
\end{figure}

The constraint is removing noise, not reshaping the partition
(Figure~\ref{fig:min-dwell}): the
STANDBY total moves by $0.00$\,h from $0$\,s to $10$\,s and by $1\%$ out
to $60$\,s, and no physics check changes verdict. Its cost is measured
directly --- $0.19\%$ of rows relabelled on this machine, $0.15$--$6.1\%$
across the plant (Section~\ref{sec:multimachine:general}) --- and it is
the component the ablation of Section~\ref{sec:ablation:state} credits
with removing flicker, in preference to the HMM.

\subsection{Validation without ground truth}
\label{sec:validation}

IMDELD supplies no per-timestamp state annotation, and obtaining one
would require instrumenting each machine's contactor. The partition
therefore cannot be scored by any supervised metric, and the alternative
adopted here is to require agreement with evidence the model does not
have access to.

Three independent bodies of evidence are used, and they fail in
different ways, which is the property that makes their agreement
informative.

\textbf{(i) Circuit physics.} The five checks C1--C5 of
Section~\ref{sec:method:accept_reject}, of which C3 (power-factor
separation) and C4 (no-load current ratio) are the two that do not
depend on power magnitude and therefore the two that can distinguish an
energised idle machine from a lightly loaded one. Pelletizer~I passes
$5/5$, with a power-factor separation of $0.72$ and a no-load current
ratio of $0.36$ --- inside the $[0.25,0.50]$ band that induction-motor
theory predicts and that no part of the fitting procedure was told
about.

\textbf{(ii) Model-level diagnostics.} A ten-test harness (T1--T10)
covering data sanity, EM convergence, model order, component separation
at $2\sigma$, assignment confidence, ghost components, physical
thresholds, variance ordering, temporal alignment with nights and
weekends, and agreement with the published plant calendar. T3 is
reported as a \emph{diagnostic} rather than a selector: as
Section~\ref{sec:multimachine:general} shows, BIC does not identify $k$
on records of this size, and the physics constraints do the selecting.

\textbf{(iii) Independently derived classifiers.} Two threshold rules
constructed from the same raw channels but with no reference to the
mixture model: a power-factor classifier and a magnetising-current-ratio
classifier, each with its threshold chosen by Otsu's
criterion~\cite{otsu1979threshold} rather than by hand, so that neither
is tuned to agree. On pelletizer~I the Otsu power-factor threshold falls
at $0.474$ and the resulting classifier agrees with the state model on
$\mathbf{99.4\%}$ of energised rows; the current-ratio classifier agrees
on $97.3\%$. Agreement across the eight machines is $85$--$99.7\%$, with
one instructive exception at $55.1\%$ on the machine whose physics
checks fail hardest (Section~\ref{sec:multimachine:states}). Verdicts
are aggregated by consensus~\cite{strehl2002cluster}.

Finally, the plant's own calendar provides an external operational
check: the facility halts every weekday between 17:00 and 22:00 local
time to avoid the peak tariff, and does not run at weekends. On the full
record $76.8\%$ of readings inside that closed window are labelled
non-productive. The converse is reported alongside it, because a
degenerate partition calling everything OFF would otherwise score
$100\%$.

\textbf{The battery is not decoration, and Phase~V measured what it is
worth.} Re-fitting the labeller at ten seeds moves the headline STANDBY
total by a coefficient of variation of $33.7\%$, and three seeds in ten
converge on a cluster that is physically a different state altogether.
The battery scores exactly those three below threshold and the other
seven at $1.00$, and conditioning on a passing score collapses the
spread to $3.0\%$ (Section~\ref{sec:significance:seeds}). That is why
the accept/reject loop is part of the method rather than an appendix to
it, and why every state-time total in this paper is quoted conditional
on it.

%  Figures are referenced by bare name; add to the preamble
%      \graphicspath{{../outputs/phase2/figures/pelletizer-I/}}
%  Bibliography: paper/references.bib
%  Numbers below are generated by
%      python run_phase2.py
%  and stored in outputs/phase2/task3/<machine>/task3_results.json,
%  outputs/phase2/task4/<machine>/{comparison_table.csv,
%  flicker_diagnosis.json,sticky_prior_sweep.csv,min_dwell_decode.csv}.
% =====================================================================

\section{The State Layer}
\label{sec:states}
\label{sec:results}

This section reports what the state-inference stage of
Section~\ref{sec:method:labelling} produces on pelletizer~I: which
measured quantities separate the states, how the choice of clustering
algorithm compares with the choice of decoding rule, and how a partition
is validated when there is nothing to validate it against.

\subsection{What separates STANDBY}
\label{sec:states:features}

Twenty-four features are derived from the six measured channels in four
groups: raw electrical (5), derived electrical (4 --- power factor,
$Q/P$ ratio, load factor, apparent-power residual), rolling statistics
(8, over a $300$\,s window), and temporal (7). Rolling windows are
\emph{trailing}, not centred, and are computed within contiguous
segments so that none spans an acquisition gap.

\textbf{The target is STANDBY against $\{$WORKING, PEAK\_LOAD$\}$ on
energised rows only.} Separating OFF from everything else is trivial and
would dominate any importance ranking computed over all four classes;
separating an energised idle machine from a lightly loaded running one
is the problem the framework actually has to solve. A four-class
ranking was also computed and is not reported here, because it measures
the OFF boundary and invites the two to be averaged.

\begin{table}[t]
\centering
\caption{Feature-group importance for the STANDBY-versus-productive
split, by three measures with different failure modes. Each column is
normalised to its own total.}
\label{tab:feature-groups}
\small
\begin{tabular}{@{}lrrr@{}}
\toprule
Group & mutual info. & split gain & SHAP \\
\midrule
raw electrical      & $0.350$ & $\mathbf{0.560}$ & $\mathbf{0.537}$ \\
derived electrical  & $0.250$ & $0.438$ & $0.310$ \\
rolling statistics  & $\mathbf{0.388}$ & $0.001$ & $0.116$ \\
temporal            & $0.012$ & $0.002$ & $0.036$ \\
\bottomrule
\end{tabular}
\end{table}

Three readings of Table~\ref{tab:feature-groups} matter for the rest of
the paper.

\textbf{The physics argument is empirically confirmed.}
\texttt{power\_factor} and \texttt{q\_p\_ratio} rank second and third by
consensus over the three measures, behind \texttt{active\_power} and
ahead of \texttt{current}, and the four derived electrical features
carry $31$--$44\%$ of total importance. Methods~1 and~2 of the
validation battery (Section~\ref{sec:validation}) were previously
justified by induction-motor theory alone; they now have a measured
basis on the machine's own data. Section~\ref{sec:explain:state}
strengthens this: under Shapley attribution power factor overtakes
active power outright.

\textbf{Temporal features are nearly irrelevant to this split}
($1$--$4\%$). That is the desirable outcome, not a null one: a state
model leaning on \texttt{hour\_of\_day} would be reproducing the shift
schedule rather than the machine's electrical condition, and would not
transfer to a facility with a different calendar. It is also the exact
opposite of what the \emph{decision} layer's forecaster does, which
draws $82\%$ of its attribution from the calendar
(Section~\ref{sec:crossmachine:decision}) --- the two stages of the
framework depend on almost disjoint information, and only the state
stage is physics-driven.

\textbf{High mutual information with near-zero split gain is the
redundancy signature.} Rolling statistics are individually informative
(MI $0.388$, the largest in the column) and add nothing once the tree
already holds the instantaneous value (gain $0.001$). Reporting either
measure alone would misdescribe them.

\paragraph{A caveat that must not be dropped.}
The surrogate used for the gain and SHAP columns reproduces its targets
to $99.9\%$ accuracy. \emph{That number is not a validation result.} The
targets are the state model's own labels and the state model was fitted
on a subset of these same features, so the surrogate is largely
re-deriving a deterministic function of its own inputs. The rankings are
meaningful; the accuracy is evidence of nothing. Validation is
Section~\ref{sec:validation}'s job.

\subsection{The emission model barely matters; the decoding rule does}
\label{sec:results:task4}

Six state-inference configurations are compared on one representation
and one protocol: $64$ contiguous two-hour blocks spread across the
record and never crossing a gap, with alternate blocks forming a fitting
set and an evaluation set of $230\,400$ rows each. DBSCAN and spectral
clustering are transductive, so their labels are extended to the
evaluation blocks by one-nearest-neighbour, applied identically to both.
Because no ground truth exists, the comparison is scored on flicker,
median dwell, the physics battery, the balanced schedule score, and the
STANDBY hour total --- never on accuracy.

\begin{table}[t]
\centering
\caption{State-inference comparison on the evaluation blocks. Flicker is
the fraction of state runs shorter than the $10$\,s physical
plausibility floor; the binary column recomputes it on the
STANDBY-versus-rest sequence the decision layer consumes. No column is
interpretable on its own --- see the spectral row.}
\label{tab:state-detection}
\small
\setlength{\tabcolsep}{4pt}
\begin{tabular}{@{}lrrrrrr@{}}
\toprule
Method & flick. & flick. & dwell & phys. & sched. & STANDBY \\
       & \%     & bin.\% & (s)   &       & \%      & (h) \\
\midrule
$k$-means \cite{lloyd1982least}      & $78.2$ & $49.1$ & $2$ & $5/5$ & $84.4$ & $4.38$ \\
DBSCAN \cite{ester1996dbscan}        & $62.5$ & $50.3$ & $4$ & $5/5$ & $84.4$ & $4.34$ \\
Spectral \cite{ng2001spectral}       & $98.2$ & $98.8$ & $2$ & $\mathbf{4/5}$ & $\mathbf{89.0}$ & $\mathbf{15.01}$ \\
GMM, i.i.d.\ \cite{mclachlan2000finite} & $66.9$ & $53.6$ & $3$ & $5/5$ & $84.5$ & $4.17$ \\
GMM--HMM \cite{rabiner1989tutorial}  & $44.5$ & $44.7$ & $13$ & $5/5$ & $84.5$ & $4.16$ \\
\textbf{GMM--HMM + min-dwell}        & $\mathbf{0.00}$ & $\mathbf{0.00}$ & $\mathbf{92.5}$ & $5/5$ & $84.5$ & $4.20$ \\
\bottomrule
\end{tabular}
\end{table}

\begin{figure*}[t]
\centering
\includegraphics[width=\textwidth]{fig_p7_state_detection}
\caption{The same six configurations as
Table~\ref{tab:state-detection}, plotted. \textbf{(A)} Flicker against
median dwell. The three filled points share one set of GMM emissions
and differ only in how those emissions are decoded; the arrows follow
i.i.d.\ assignment $\rightarrow$ Viterbi $\rightarrow$ Viterbi under a
minimum-dwell floor. The three squares are alternative emission models
decoded i.i.d.\ Changing the emission model moves the result within one
cluster; changing the decoding rule moves it two orders of magnitude in
dwell and to zero in flicker. \textbf{(B)} Inferred STANDBY hours.
Spectral clustering assigns $3.6\times$ every other method's total while
scoring the highest factory-schedule agreement of any row in the table,
which is why the battery is argued for in preference to any single
ground-truth-free metric. Hatching marks a failed physics check in both
panels.}
\label{fig:state-detection}
\end{figure*}

\textbf{Choosing a better clustering algorithm buys almost nothing;
choosing a better sequence model buys everything} ---
Figure~\ref{fig:state-detection}(A) is the table read as a picture, and
the two families separate along different axes. $k$-means, DBSCAN and
the GMM all pass $5/5$ physics checks, agree with the factory schedule
to within $0.2$ points, and find $4.17$--$4.38$ STANDBY hours on the
evaluation blocks. Switching from i.i.d.\ assignment to Viterbi decoding
over the \emph{same} emissions cuts flicker from $66.9\%$ to $44.5\%$
and quadruples the median dwell; adding the minimum-dwell constraint
takes flicker to zero and the median dwell to $92.5$\,s. The
emission-model rows span $62$--$78\%$ of flicker; the decoding rows span
$0$--$67\%$.

The comparison is clean in the one place it needs to be: the
\textsc{gmm} and \textsc{gmm--hmm} rows differ \emph{only} in the
decoding rule, because the HMM's emissions are warm-started from the
fitted GMM and held fixed (Section~\ref{sec:method:ghmm}). Fitting HMM
emissions from scratch by Baum--Welch converges to a materially worse
optimum on this record --- its ``STANDBY'' absorbs part of the
productive range at $10.3$\,kW instead of finding the no-load band at
$3.67$\,kW, and it fails the variance-ordering check --- which would
have confounded the value of temporal decoding with EM initialisation
luck.

\textbf{Spectral clustering is the one genuine failure, and it is
instructive} (Figure~\ref{fig:state-detection}(B)). It assigns $15.01$
STANDBY hours, $3.6\times$ every other
method, with $98.8\%$ binary flicker and a no-load current ratio of
$0.001$ against the acceptance band $[0.25,0.50]$. It also scores the
\emph{highest} schedule agreement in the table ($89.0\%$): a partition
that over-assigns non-productive states automatically looks good on the
factory-closed window. That is a weakness of the schedule metric in
isolation and the reason the STANDBY hour total is printed beside it.
\textbf{No single ground-truth-free metric can be quoted alone}, which
is the argument for the battery rather than for a favourite score.

\paragraph{A protocol correction worth recording.}
DBSCAN chooses its own cluster count and returned $28$. Naming the four
lowest-power clusters with the four canonical state names put all four
``states'' inside the machine's low-power range and produced an apparent
total failure --- $0.06$ STANDBY hours, $98.9\%$ flicker --- that was
entirely an artefact of the mapping. Cluster means are now grouped into
four size-weighted power bands, so sub-clusters of one physical state
are re-merged, and DBSCAN passes $5/5$ with $4.34$ STANDBY hours. Any
comparison against a method that over-segments needs this step, and
without it this paper would have reported a baseline failure that does
not exist.

\subsection{Flicker: why a transition prior cannot remove it}
\label{sec:flicker}

The expectation going in was that Viterbi decoding under a sticky prior
would take flicker to a few per cent. It does not, and the reason is
arithmetic rather than empirical.

\textbf{The sticky prior as conventionally set is numerically inert.}
Setting the Dirichlet pseudo-count to the target dwell in samples gives
$\kappa=30$ against roughly $576\,000$ observed transitions. A flat
prior and $\kappa=30$ give flicker $64.10\%$ and $64.09\%$ and minimum
self-transition $0.9677$ and $0.9678$: the prior is present in the
model, cited in its documentation, and has no measurable effect.

\textbf{Strengthening it does not help either.} Scaling the prior to ten
times the total evidence ($\kappa=575\,920$, minimum self-transition
$0.9994$, implied mean dwell $1\,605$\,s) still leaves flicker at
$50.2\%$.

\textbf{The bound is structural.} In a Viterbi decode the cost of
changing state is $\log(p_{\text{self}}/p_{\text{switch}})$, measured
here at $9.23$ nats at its largest and bounded above however strong the
prior is made. The emission log-likelihood gap between the best and
second-best state has median $27.8$ nats and $90$th percentile
$41\,701$ nats. \textbf{At $55.7\%$ of timesteps the emission term
exceeds the largest penalty the transition matrix can apply}, so the
transition model is arithmetically incapable of changing the decision
there and the HMM degenerates toward per-sample assignment. This is a
general property of applying an HMM to low-noise, multi-channel
measurements, not a property of this record.

\textbf{What fixes it is an explicit duration constraint} --- the
discrete analogue of a hidden semi-Markov duration
model~\cite{yu2010hsmm}, imposed at decode time. The floor is set equal
to the flicker threshold itself, so it removes exactly what the metric
defines as implausible and nothing more.
Table~\ref{tab:min-dwell} sweeps the floor over an order of magnitude.

\begin{table}[t]
\centering
\caption{Minimum-dwell sweep. Flicker goes to zero at no cost to the
physics checks, the schedule agreement, or the STANDBY total. This
sweep is run on the flicker-diagnosis decode rather than on the
evaluation blocks of Table~\ref{tab:state-detection}, which is why its
absolute physics score and STANDBY total differ from that table's; the
comparison of interest is down the column, not across the two tables.}
\label{tab:min-dwell}
\small
\begin{tabular}{@{}rrrrr@{}}
\toprule
min.\ dwell & flicker \% & median dwell (s) & physics & STANDBY (h) \\
\midrule
$0$\,s  & $64.09$ & $4$   & $4/5$ & $5.34$ \\
$10$\,s & $\mathbf{0.00}$ & $153$ & $4/5$ & $5.34$ \\
$30$\,s & $0.00$ & $327.5$ & $4/5$ & $5.37$ \\
$60$\,s & $0.00$ & $436.5$ & $4/5$ & $5.40$ \\
\bottomrule
\end{tabular}
\end{table}

\begin{figure}[t]
\centering
\includegraphics[width=\columnwidth]{fig_p7_min_dwell}
\caption{The minimum-dwell sweep of Table~\ref{tab:min-dwell}.
\textbf{(A)} Flicker under both definitions collapses to $0.00\%$ at the
first non-zero floor and stays there. \textbf{(B)} The median dwell
rises by two orders of magnitude while the inferred STANDBY total is
flat. The right-hand axis is anchored at zero deliberately: on an
auto-scaled axis the $1\%$ drift across the whole sweep would render as
a pronounced trend and say the opposite of what the data says.}
\label{fig:min-dwell}
\end{figure}

The constraint is removing noise, not reshaping the partition
(Figure~\ref{fig:min-dwell}): the
STANDBY total moves by $0.00$\,h from $0$\,s to $10$\,s and by $1\%$ out
to $60$\,s, and no physics check changes verdict. Its cost is measured
directly --- $0.19\%$ of rows relabelled on this machine, $0.15$--$6.1\%$
across the plant (Section~\ref{sec:multimachine:general}) --- and it is
the component the ablation of Section~\ref{sec:ablation:state} credits
with removing flicker, in preference to the HMM.

\subsection{Validation without ground truth}
\label{sec:validation}

IMDELD supplies no per-timestamp state annotation, and obtaining one
would require instrumenting each machine's contactor. The partition
therefore cannot be scored by any supervised metric, and the alternative
adopted here is to require agreement with evidence the model does not
have access to.

Three independent bodies of evidence are used, and they fail in
different ways, which is the property that makes their agreement
informative.

\textbf{(i) Circuit physics.} The five checks C1--C5 of
Section~\ref{sec:method:accept_reject}, of which C3 (power-factor
separation) and C4 (no-load current ratio) are the two that do not
depend on power magnitude and therefore the two that can distinguish an
energised idle machine from a lightly loaded one. Pelletizer~I passes
$5/5$, with a power-factor separation of $0.72$ and a no-load current
ratio of $0.36$ --- inside the $[0.25,0.50]$ band that induction-motor
theory predicts and that no part of the fitting procedure was told
about.

\textbf{(ii) Model-level diagnostics.} A ten-test harness (T1--T10)
covering data sanity, EM convergence, model order, component separation
at $2\sigma$, assignment confidence, ghost components, physical
thresholds, variance ordering, temporal alignment with nights and
weekends, and agreement with the published plant calendar. T3 is
reported as a \emph{diagnostic} rather than a selector: as
Section~\ref{sec:multimachine:general} shows, BIC does not identify $k$
on records of this size, and the physics constraints do the selecting.

\textbf{(iii) Independently derived classifiers.} Two threshold rules
constructed from the same raw channels but with no reference to the
mixture model: a power-factor classifier and a magnetising-current-ratio
classifier, each with its threshold chosen by Otsu's
criterion~\cite{otsu1979threshold} rather than by hand, so that neither
is tuned to agree. On pelletizer~I the Otsu power-factor threshold falls
at $0.474$ and the resulting classifier agrees with the state model on
$\mathbf{99.4\%}$ of energised rows; the current-ratio classifier agrees
on $97.3\%$. Agreement across the eight machines is $85$--$99.7\%$, with
one instructive exception at $55.1\%$ on the machine whose physics
checks fail hardest (Section~\ref{sec:multimachine:states}). Verdicts
are aggregated by consensus~\cite{strehl2002cluster}.

Finally, the plant's own calendar provides an external operational
check: the facility halts every weekday between 17:00 and 22:00 local
time to avoid the peak tariff, and does not run at weekends. On the full
record $76.8\%$ of readings inside that closed window are labelled
non-productive. The converse is reported alongside it, because a
degenerate partition calling everything OFF would otherwise score
$100\%$.

\textbf{The battery is not decoration, and Phase~V measured what it is
worth.} Re-fitting the labeller at ten seeds moves the headline STANDBY
total by a coefficient of variation of $33.7\%$, and three seeds in ten
converge on a cluster that is physically a different state altogether.
The battery scores exactly those three below threshold and the other
seven at $1.00$, and conditioning on a passing score collapses the
spread to $3.0\%$ (Section~\ref{sec:significance:seeds}). That is why
the accept/reject loop is part of the method rather than an appendix to
it, and why every state-time total in this paper is quoted conditional
on it.

%  Figures are referenced by bare name; add to the preamble
%      \graphicspath{{../outputs/phase2/figures/pelletizer-I/}}
%  Bibliography: paper/references.bib
%  Numbers below are generated by
%      python run_phase2.py
%  and stored in outputs/phase2/task3/<machine>/task3_results.json,
%  outputs/phase2/task4/<machine>/{comparison_table.csv,
%  flicker_diagnosis.json,sticky_prior_sweep.csv,min_dwell_decode.csv}.
% =====================================================================

\section{The State Layer}
\label{sec:states}
\label{sec:results}

This section reports what the state-inference stage of
Section~\ref{sec:method:labelling} produces on pelletizer~I: which
measured quantities separate the states, how the choice of clustering
algorithm compares with the choice of decoding rule, and how a partition
is validated when there is nothing to validate it against.

\subsection{What separates STANDBY}
\label{sec:states:features}

Twenty-four features are derived from the six measured channels in four
groups: raw electrical (5), derived electrical (4 --- power factor,
$Q/P$ ratio, load factor, apparent-power residual), rolling statistics
(8, over a $300$\,s window), and temporal (7). Rolling windows are
\emph{trailing}, not centred, and are computed within contiguous
segments so that none spans an acquisition gap.

\textbf{The target is STANDBY against $\{$WORKING, PEAK\_LOAD$\}$ on
energised rows only.} Separating OFF from everything else is trivial and
would dominate any importance ranking computed over all four classes;
separating an energised idle machine from a lightly loaded running one
is the problem the framework actually has to solve. A four-class
ranking was also computed and is not reported here, because it measures
the OFF boundary and invites the two to be averaged.

\begin{table}[t]
\centering
\caption{Feature-group importance for the STANDBY-versus-productive
split, by three measures with different failure modes. Each column is
normalised to its own total.}
\label{tab:feature-groups}
\small
\begin{tabular}{@{}lrrr@{}}
\toprule
Group & mutual info. & split gain & SHAP \\
\midrule
raw electrical      & $0.350$ & $\mathbf{0.560}$ & $\mathbf{0.537}$ \\
derived electrical  & $0.250$ & $0.438$ & $0.310$ \\
rolling statistics  & $\mathbf{0.388}$ & $0.001$ & $0.116$ \\
temporal            & $0.012$ & $0.002$ & $0.036$ \\
\bottomrule
\end{tabular}
\end{table}

Three readings of Table~\ref{tab:feature-groups} matter for the rest of
the paper.

\textbf{The physics argument is empirically confirmed.}
\texttt{power\_factor} and \texttt{q\_p\_ratio} rank second and third by
consensus over the three measures, behind \texttt{active\_power} and
ahead of \texttt{current}, and the four derived electrical features
carry $31$--$44\%$ of total importance. Methods~1 and~2 of the
validation battery (Section~\ref{sec:validation}) were previously
justified by induction-motor theory alone; they now have a measured
basis on the machine's own data. Section~\ref{sec:explain:state}
strengthens this: under Shapley attribution power factor overtakes
active power outright.

\textbf{Temporal features are nearly irrelevant to this split}
($1$--$4\%$). That is the desirable outcome, not a null one: a state
model leaning on \texttt{hour\_of\_day} would be reproducing the shift
schedule rather than the machine's electrical condition, and would not
transfer to a facility with a different calendar. It is also the exact
opposite of what the \emph{decision} layer's forecaster does, which
draws $82\%$ of its attribution from the calendar
(Section~\ref{sec:crossmachine:decision}) --- the two stages of the
framework depend on almost disjoint information, and only the state
stage is physics-driven.

\textbf{High mutual information with near-zero split gain is the
redundancy signature.} Rolling statistics are individually informative
(MI $0.388$, the largest in the column) and add nothing once the tree
already holds the instantaneous value (gain $0.001$). Reporting either
measure alone would misdescribe them.

\paragraph{A caveat that must not be dropped.}
The surrogate used for the gain and SHAP columns reproduces its targets
to $99.9\%$ accuracy. \emph{That number is not a validation result.} The
targets are the state model's own labels and the state model was fitted
on a subset of these same features, so the surrogate is largely
re-deriving a deterministic function of its own inputs. The rankings are
meaningful; the accuracy is evidence of nothing. Validation is
Section~\ref{sec:validation}'s job.

\subsection{The emission model barely matters; the decoding rule does}
\label{sec:results:task4}

Six state-inference configurations are compared on one representation
and one protocol: $64$ contiguous two-hour blocks spread across the
record and never crossing a gap, with alternate blocks forming a fitting
set and an evaluation set of $230\,400$ rows each. DBSCAN and spectral
clustering are transductive, so their labels are extended to the
evaluation blocks by one-nearest-neighbour, applied identically to both.
Because no ground truth exists, the comparison is scored on flicker,
median dwell, the physics battery, the balanced schedule score, and the
STANDBY hour total --- never on accuracy.

\begin{table}[t]
\centering
\caption{State-inference comparison on the evaluation blocks. Flicker is
the fraction of state runs shorter than the $10$\,s physical
plausibility floor; the binary column recomputes it on the
STANDBY-versus-rest sequence the decision layer consumes. No column is
interpretable on its own --- see the spectral row.}
\label{tab:state-detection}
\small
\setlength{\tabcolsep}{4pt}
\begin{tabular}{@{}lrrrrrr@{}}
\toprule
Method & flick. & flick. & dwell & phys. & sched. & STANDBY \\
       & \%     & bin.\% & (s)   &       & \%      & (h) \\
\midrule
$k$-means \cite{lloyd1982least}      & $78.2$ & $49.1$ & $2$ & $5/5$ & $84.4$ & $4.38$ \\
DBSCAN \cite{ester1996dbscan}        & $62.5$ & $50.3$ & $4$ & $5/5$ & $84.4$ & $4.34$ \\
Spectral \cite{ng2001spectral}       & $98.2$ & $98.8$ & $2$ & $\mathbf{4/5}$ & $\mathbf{89.0}$ & $\mathbf{15.01}$ \\
GMM, i.i.d.\ \cite{mclachlan2000finite} & $66.9$ & $53.6$ & $3$ & $5/5$ & $84.5$ & $4.17$ \\
GMM--HMM \cite{rabiner1989tutorial}  & $44.5$ & $44.7$ & $13$ & $5/5$ & $84.5$ & $4.16$ \\
\textbf{GMM--HMM + min-dwell}        & $\mathbf{0.00}$ & $\mathbf{0.00}$ & $\mathbf{92.5}$ & $5/5$ & $84.5$ & $4.20$ \\
\bottomrule
\end{tabular}
\end{table}

\begin{figure*}[t]
\centering
\includegraphics[width=\textwidth]{fig_p7_state_detection}
\caption{The same six configurations as
Table~\ref{tab:state-detection}, plotted. \textbf{(A)} Flicker against
median dwell. The three filled points share one set of GMM emissions
and differ only in how those emissions are decoded; the arrows follow
i.i.d.\ assignment $\rightarrow$ Viterbi $\rightarrow$ Viterbi under a
minimum-dwell floor. The three squares are alternative emission models
decoded i.i.d.\ Changing the emission model moves the result within one
cluster; changing the decoding rule moves it two orders of magnitude in
dwell and to zero in flicker. \textbf{(B)} Inferred STANDBY hours.
Spectral clustering assigns $3.6\times$ every other method's total while
scoring the highest factory-schedule agreement of any row in the table,
which is why the battery is argued for in preference to any single
ground-truth-free metric. Hatching marks a failed physics check in both
panels.}
\label{fig:state-detection}
\end{figure*}

\textbf{Choosing a better clustering algorithm buys almost nothing;
choosing a better sequence model buys everything} ---
Figure~\ref{fig:state-detection}(A) is the table read as a picture, and
the two families separate along different axes. $k$-means, DBSCAN and
the GMM all pass $5/5$ physics checks, agree with the factory schedule
to within $0.2$ points, and find $4.17$--$4.38$ STANDBY hours on the
evaluation blocks. Switching from i.i.d.\ assignment to Viterbi decoding
over the \emph{same} emissions cuts flicker from $66.9\%$ to $44.5\%$
and quadruples the median dwell; adding the minimum-dwell constraint
takes flicker to zero and the median dwell to $92.5$\,s. The
emission-model rows span $62$--$78\%$ of flicker; the decoding rows span
$0$--$67\%$.

The comparison is clean in the one place it needs to be: the
\textsc{gmm} and \textsc{gmm--hmm} rows differ \emph{only} in the
decoding rule, because the HMM's emissions are warm-started from the
fitted GMM and held fixed (Section~\ref{sec:method:ghmm}). Fitting HMM
emissions from scratch by Baum--Welch converges to a materially worse
optimum on this record --- its ``STANDBY'' absorbs part of the
productive range at $10.3$\,kW instead of finding the no-load band at
$3.67$\,kW, and it fails the variance-ordering check --- which would
have confounded the value of temporal decoding with EM initialisation
luck.

\textbf{Spectral clustering is the one genuine failure, and it is
instructive} (Figure~\ref{fig:state-detection}(B)). It assigns $15.01$
STANDBY hours, $3.6\times$ every other
method, with $98.8\%$ binary flicker and a no-load current ratio of
$0.001$ against the acceptance band $[0.25,0.50]$. It also scores the
\emph{highest} schedule agreement in the table ($89.0\%$): a partition
that over-assigns non-productive states automatically looks good on the
factory-closed window. That is a weakness of the schedule metric in
isolation and the reason the STANDBY hour total is printed beside it.
\textbf{No single ground-truth-free metric can be quoted alone}, which
is the argument for the battery rather than for a favourite score.

\paragraph{A protocol correction worth recording.}
DBSCAN chooses its own cluster count and returned $28$. Naming the four
lowest-power clusters with the four canonical state names put all four
``states'' inside the machine's low-power range and produced an apparent
total failure --- $0.06$ STANDBY hours, $98.9\%$ flicker --- that was
entirely an artefact of the mapping. Cluster means are now grouped into
four size-weighted power bands, so sub-clusters of one physical state
are re-merged, and DBSCAN passes $5/5$ with $4.34$ STANDBY hours. Any
comparison against a method that over-segments needs this step, and
without it this paper would have reported a baseline failure that does
not exist.

\subsection{Flicker: why a transition prior cannot remove it}
\label{sec:flicker}

The expectation going in was that Viterbi decoding under a sticky prior
would take flicker to a few per cent. It does not, and the reason is
arithmetic rather than empirical.

\textbf{The sticky prior as conventionally set is numerically inert.}
Setting the Dirichlet pseudo-count to the target dwell in samples gives
$\kappa=30$ against roughly $576\,000$ observed transitions. A flat
prior and $\kappa=30$ give flicker $64.10\%$ and $64.09\%$ and minimum
self-transition $0.9677$ and $0.9678$: the prior is present in the
model, cited in its documentation, and has no measurable effect.

\textbf{Strengthening it does not help either.} Scaling the prior to ten
times the total evidence ($\kappa=575\,920$, minimum self-transition
$0.9994$, implied mean dwell $1\,605$\,s) still leaves flicker at
$50.2\%$.

\textbf{The bound is structural.} In a Viterbi decode the cost of
changing state is $\log(p_{\text{self}}/p_{\text{switch}})$, measured
here at $9.23$ nats at its largest and bounded above however strong the
prior is made. The emission log-likelihood gap between the best and
second-best state has median $27.8$ nats and $90$th percentile
$41\,701$ nats. \textbf{At $55.7\%$ of timesteps the emission term
exceeds the largest penalty the transition matrix can apply}, so the
transition model is arithmetically incapable of changing the decision
there and the HMM degenerates toward per-sample assignment. This is a
general property of applying an HMM to low-noise, multi-channel
measurements, not a property of this record.

\textbf{What fixes it is an explicit duration constraint} --- the
discrete analogue of a hidden semi-Markov duration
model~\cite{yu2010hsmm}, imposed at decode time. The floor is set equal
to the flicker threshold itself, so it removes exactly what the metric
defines as implausible and nothing more.
Table~\ref{tab:min-dwell} sweeps the floor over an order of magnitude.

\begin{table}[t]
\centering
\caption{Minimum-dwell sweep. Flicker goes to zero at no cost to the
physics checks, the schedule agreement, or the STANDBY total. This
sweep is run on the flicker-diagnosis decode rather than on the
evaluation blocks of Table~\ref{tab:state-detection}, which is why its
absolute physics score and STANDBY total differ from that table's; the
comparison of interest is down the column, not across the two tables.}
\label{tab:min-dwell}
\small
\begin{tabular}{@{}rrrrr@{}}
\toprule
min.\ dwell & flicker \% & median dwell (s) & physics & STANDBY (h) \\
\midrule
$0$\,s  & $64.09$ & $4$   & $4/5$ & $5.34$ \\
$10$\,s & $\mathbf{0.00}$ & $153$ & $4/5$ & $5.34$ \\
$30$\,s & $0.00$ & $327.5$ & $4/5$ & $5.37$ \\
$60$\,s & $0.00$ & $436.5$ & $4/5$ & $5.40$ \\
\bottomrule
\end{tabular}
\end{table}

\begin{figure}[t]
\centering
\includegraphics[width=\columnwidth]{fig_p7_min_dwell}
\caption{The minimum-dwell sweep of Table~\ref{tab:min-dwell}.
\textbf{(A)} Flicker under both definitions collapses to $0.00\%$ at the
first non-zero floor and stays there. \textbf{(B)} The median dwell
rises by two orders of magnitude while the inferred STANDBY total is
flat. The right-hand axis is anchored at zero deliberately: on an
auto-scaled axis the $1\%$ drift across the whole sweep would render as
a pronounced trend and say the opposite of what the data says.}
\label{fig:min-dwell}
\end{figure}

The constraint is removing noise, not reshaping the partition
(Figure~\ref{fig:min-dwell}): the
STANDBY total moves by $0.00$\,h from $0$\,s to $10$\,s and by $1\%$ out
to $60$\,s, and no physics check changes verdict. Its cost is measured
directly --- $0.19\%$ of rows relabelled on this machine, $0.15$--$6.1\%$
across the plant (Section~\ref{sec:multimachine:general}) --- and it is
the component the ablation of Section~\ref{sec:ablation:state} credits
with removing flicker, in preference to the HMM.

\subsection{Validation without ground truth}
\label{sec:validation}

IMDELD supplies no per-timestamp state annotation, and obtaining one
would require instrumenting each machine's contactor. The partition
therefore cannot be scored by any supervised metric, and the alternative
adopted here is to require agreement with evidence the model does not
have access to.

Three independent bodies of evidence are used, and they fail in
different ways, which is the property that makes their agreement
informative.

\textbf{(i) Circuit physics.} The five checks C1--C5 of
Section~\ref{sec:method:accept_reject}, of which C3 (power-factor
separation) and C4 (no-load current ratio) are the two that do not
depend on power magnitude and therefore the two that can distinguish an
energised idle machine from a lightly loaded one. Pelletizer~I passes
$5/5$, with a power-factor separation of $0.72$ and a no-load current
ratio of $0.36$ --- inside the $[0.25,0.50]$ band that induction-motor
theory predicts and that no part of the fitting procedure was told
about.

\textbf{(ii) Model-level diagnostics.} A ten-test harness (T1--T10)
covering data sanity, EM convergence, model order, component separation
at $2\sigma$, assignment confidence, ghost components, physical
thresholds, variance ordering, temporal alignment with nights and
weekends, and agreement with the published plant calendar. T3 is
reported as a \emph{diagnostic} rather than a selector: as
Section~\ref{sec:multimachine:general} shows, BIC does not identify $k$
on records of this size, and the physics constraints do the selecting.

\textbf{(iii) Independently derived classifiers.} Two threshold rules
constructed from the same raw channels but with no reference to the
mixture model: a power-factor classifier and a magnetising-current-ratio
classifier, each with its threshold chosen by Otsu's
criterion~\cite{otsu1979threshold} rather than by hand, so that neither
is tuned to agree. On pelletizer~I the Otsu power-factor threshold falls
at $0.474$ and the resulting classifier agrees with the state model on
$\mathbf{99.4\%}$ of energised rows; the current-ratio classifier agrees
on $97.3\%$. Agreement across the eight machines is $85$--$99.7\%$, with
one instructive exception at $55.1\%$ on the machine whose physics
checks fail hardest (Section~\ref{sec:multimachine:states}). Verdicts
are aggregated by consensus~\cite{strehl2002cluster}.

Finally, the plant's own calendar provides an external operational
check: the facility halts every weekday between 17:00 and 22:00 local
time to avoid the peak tariff, and does not run at weekends. On the full
record $76.8\%$ of readings inside that closed window are labelled
non-productive. The converse is reported alongside it, because a
degenerate partition calling everything OFF would otherwise score
$100\%$.

\textbf{The battery is not decoration, and Phase~V measured what it is
worth.} Re-fitting the labeller at ten seeds moves the headline STANDBY
total by a coefficient of variation of $33.7\%$, and three seeds in ten
converge on a cluster that is physically a different state altogether.
The battery scores exactly those three below threshold and the other
seven at $1.00$, and conditioning on a passing score collapses the
spread to $3.0\%$ (Section~\ref{sec:significance:seeds}). That is why
the accept/reject loop is part of the method rather than an appendix to
it, and why every state-time total in this paper is quoted conditional
on it.
